% !TeX root = asymptotic-behavior-of-divisors.tex
Let $\kkk$ be an algebraically closed field, usually of characteristic $0$.
    Let $X$ be a normal projective variety over $\kkk$ of dimension $d$.
    
\subsection{Volume and bigness}

    \begin{definition}
        Let $D$ be a divisor on a projective variety $X$ of dimension $d$.
        The \emph{volume} of $D$ is defined as
        \[ \vol(D) \coloneqq \lim_{m \to +\infty} \frac{h^0(X, \calO_X(mD))}{m^d/d!}. \]
    \end{definition}


    \begin{theorem}
        Let $D$ be a divisor on a projective variety $X$ of dimension $d$.
        \begin{enumerate}
            \item The limit in the definition of $\vol(D)$ exists.
            \item If $D$ is ample, then $\vol(D) = D^d$.
            \item The volume $\vol(D)$ depends only on the numerical equivalence class of $D$.
            \item The volume function $\vol: \NS(X)_\bbQ \to \bbR$ is continuous and homogeneous of degree $d$.
        \end{enumerate}
    \end{theorem}


    \begin{definition}
        A divisor $D$ on a projective variety $X$ is called \emph{big} if $\vol(D) > 0$.
    \end{definition}


    \begin{theorem}
        Let $D$ be a divisor on a projective variety $X$ of dimension $d$.
        \begin{enumerate}
            \item If $D \equiv D'$ and $D$ is big, then $D'$ is also big.
            \item If $D$ is big, then there exist an ample divisor $A$ and an effective divisor $E$ such that $D \sim_\bbQ A + E$.
            \item A numerical class $D \in \NS(X)_\bbQ$ is big if and only if it lies in the interior of the pseudo-effective cone $\Psef^1(X)$.
        \end{enumerate}
    \end{theorem}

    \begin{theorem}[{Siu's criterion, ref. \cite[Theorem 2.2.16]{Laz04a}}]
        Let $D$ be a divisor on a projective variety $X$ of dimension $d$.
        Then $D$ is big if and only if there exist an ample divisor $A$ and an effective divisor $E$ such that $D \sim_\bbQ A + E$.
        \Yang{}
    \end{theorem}

    \begin{corollary}
        If $D$ is a nef divisor on a projective variety $X$ of dimension $d$ such that $D^d > 0$, then $D$ is big.
    \end{corollary}

    \begin{theorem}
        Volume is log-concave on the big cone, i.e. for any big divisors $D_1, D_2$ on a projective variety $X$ of dimension $d$, we have
        \[ \vol(D_1 + D_2)^{1/d} \geq \vol(D_1)^{1/d} + \vol(D_2)^{1/d}. \]
    \end{theorem}

    \begin{theorem}
        Volume is semi-continuous on a family.
    \end{theorem}


    

% \subsection{Zariski decomposition}


% \subsection{BDPP's criterion}

%     \begin{theorem}[{BDPP's criterion, ref. \cite[Theorem 0.2]{BDPP12}}]
%         Let $X$ be a smooth projective variety over $\kkk$.
%         Recall that $\Nef^{d-1}(X) \subset \upN^{d-1}(X)$ is defined as the dual cone of $\Psef_{d-1}(X) \subset \upN_{d-1}(X)$.
%         Then we have the following equivalence of conditions describing $\Nef^{d-1}(X)$:
%         \begin{enumerate}
%             \item $\Nef^{d-1}(X) = \Psef^1(X)^\vee$.
%             \item $\Nef^{d-1}(X) = \overline{\Mov}_1(X)$.
%             \item $\Nef^{d-1}(X)$ is the closure of the cone generated by classes of form $p_*(H_1 \cdots H_{d-1})$ where $p: Y \to X$ is a birational morphism and $H_i$ are nef divisors on $Y$.
%         \end{enumerate}
%         \Yang{}
%     \end{theorem}


\subsection{Base locus}

    \begin{example}
        Let $E$ be an elliptic curve and $P \in E(\kkk)$.
        Set $X = \bbP(\calO_E \oplus \calO_E(-P))$ with the natural projection $\pi: X \to E$ and the unique minimal section $C_0$ with $C_0^2 = -1$.
        Let $D = C_0 + \pi^*P$.
        Then 
        \[ \Bs |D| = \{Q\}, \]
        where $Q$ is some point on $X$ projecting to $P$.

        Indeed, we have 
        \[ h^0(X, D) = h^0(E, \pi_*\left(\calO_X(C_0) \otimes \pi^*\calO_E(P) \right)) = h^0(E, \calO_E \oplus \calO_E(P)) = 2. \]
        Thus $\dim |D| = 1$.
        For $D_t \in |D|$, write $D_t = H_t + \pi^*B$ where $H_t$ is the horizontal part and $B \in \Div(E)$ is the vertical part.
        Intersecting with fibers, we get that $H_t$ is prime and a section of $\pi$.
        By Bernstein's theorem, for general $t$, $D_t$ is irreducible and thus $D_t = H_t$.
        Since $D^2 = 1$, we only need to find the intersection of two distinct members of $|D|$ to find the base locus.
        Note on $\CH_0(X)$, we have 
        \[ D_t \cap D_s = (C_0 + \pi^*P) \cdot (C_0 + \pi^*P) = [Q] \]
        for a point $Q$.
        By proper pushforward of $\CH_0$, we see that $Q$ projects to $P$.
        It follows that any two irreducible members of $|D|$ intersect at $Q$ since they are section and $\Bs |D| \subset \{Q\}$.
        
        Fix an irreducible member $D_0 \in |D|$.
        For every $D_t = H_t + \pi^*B_t \in |D|$, $H_t$ is different from $D_0$ and thus 
        \[ D_t \cap D_0 = [H_t \cap D_0] + [\pi^*B_t \cap D_0] = [Q]. \]
        Note both $[H_t \cap D_0]$ and $[\pi^*B_t \cap D_0]$ are effective, we get that $[\pi^*B_t \cap D_0] = [Q]$ or $0$.
        If $[\pi^*B_t \cap D_0] = [Q]$, then $B_t = P$ and thus $D_t = H_t + \pi^*P \sim C_0 + \pi^*P$.
        Then $H_t \sim C_0$ and thus $H_t = C_0$.
        It pass through $Q$.
        Otherwise, $[\pi^*B_t \cap D_0] = 0$ and thus $B_t = 0$ and thus $D_t$ is irreducible.
        We have seen that $D_t$ also pass through $Q$.
        Hence, $\Bs |D| = \{Q\}$.

        % Hence $Q$ is the unique base point of $|D|$.
    \end{example}
    
